\section{Symmetric subspaces}

% \begin{lemma}\label{TestLemma}
%   Test lemma
% \end{lemma}

% \begin{theorem}\uses{TestLemma}\label{TestTheorem}
%   Test theorem
% \end{theorem}

% Yan: everything we define as \begin{definition, proposition, theorem, lemma} will show up in the
% dependency graph, so we can use the \uses{} command to specify the dependencies between them.
Suppose $\mathcal{X} = \mathbb{C}^{[d]}$ and $\mathcal{Y} = \mathbb{C}^{[d]}$ where $d \in \mathbb{N}$.
Let the set of linear operators be $\mathcal{L}(\mathcal{X}, \mathcal{Y})$.
When $\mathcal{Y} = \mathcal{X}$, let the set of linear operators be $\mathcal{L}(\mathcal{X})$.

\begin{definition}{standard computational basis of linear operators}\label{def:standard_basis_Lk}
  The {standard computational basis of linear operators} $\{E_{i, j}\}_{i, j \in [d]}$ is defined as the set of $d^2$ operators in $\mathcal{L}(\mathcal{X})$ such that
\begin{equation}
    E_{i, j} = |i \rangle \langle j|
\end{equation}
for all $i, j \in [d]$, where $\{|i\rangle\}_{i \in [d]}$ is the standard computational basis of $\mathcal{X}$.
\end{definition}

\begin{definition}{operator-vector correspondence}\label{def:operator_vector_correspondence}
    The {operator-vector correspondence} is a bijective mapping between the set of linear operators $\mathcal{L}(\mathcal{X})$ and the tensor product space $\mathcal{X} \otimes \mathcal{X}$ defined as
    \begin{equation}
    \text{vec} : \mathcal{L}(\mathcal{X}) \to \mathcal{X} \otimes \mathcal{X}, \quad E_{i, j} \mapsto |i\rangle \otimes |j\rangle\text{ for all } i, j \in [d].
    \end{equation}

\end{definition}

\begin{definition}[symmetric subspace of linear operators]\label{def:sym_Lk}
  The {symmetric subspace of linear operators} $\mathcal{L}(\mathcal{X})^{\varovee k}$ is the set of all operators $X \in \mathcal{L}(\mathcal{X}^{\otimes k})$ that remain unchanged (invariant) under the permutation $W_{\pi}$.
That is,
\begin{equation}
    \mathcal{L}(\mathcal{X})^{\varovee k} = \{X \in \mathcal{L}(\mathcal{X}^{\otimes k}) : X = W_{\pi} X W_{\pi}^{*} \text{ for all } \pi \in S_{k} \}.
    \label{eq:sym_Lk}
\end{equation}

\end{definition}

% Yan: If we want to add certain definitions that probably already exists in mathlib,
% we can add them here as a remark, and we can reference them for the reader's convenience.
\begin{remark}
Let $S_k$ denote the symmetric group. Then given $\pi \in S_k$, the permutation matrix (operator) $V_d(\pi)$ is the unitary matrix that satisfies
\begin{equation}\label{eq:permutation}
    W_{\pi} |\psi_1 \rangle \otimes \dots |\psi_k = |\psi_{\pi^{-1}(1)}\rangle \otimes \dots \otimes |\psi_{\pi^{-1}(k)}\rangle
\end{equation}
for all $|\psi_1 \rangle, \dots, |\psi_k\rangle \in \mathbb{C}^{d}$.
\end{remark}

\begin{remark}\label{remark:LieBracket}
For any pair of operators $X, Y \in \mathcal{L}(\mathcal{X})$, the Lie bracket $[X, Y] \in \mathcal{L}(\mathcal{X})$ is defined as
\begin{equation}
    [X, Y] = XY - YX,
\end{equation}
such that $[X, Y] = 0$ if and only if $X$ and $Y$ commute.

\end{remark}

\begin{remark}\uses{remark:LieBracket}\label{remark:Commutant}
Let $\mathcal{A} \subseteq \mathcal{L}(\mathcal{X})$ be a subalgebra of $\mathcal{L}(\mathcal{X})$.
The commutant of $\mathcal{A}$ is given by
\begin{equation}
    \text{Comm}(\mathcal{A}) = \{Y \in \mathcal{L}(\mathcal{X}) : [X, Y] = 0 \text{ for all } X \in \mathcal{A}\}
\end{equation}

\end{remark}

\begin{definition}{diagonal operators}\label{def:diagonal_operators}
An operator $X \in \mathcal{L}(\mathcal{X})$ is a diagonal operator if $X_{a, b} = 0$ for all $a, b \in [d]$ where $a \neq b$.
For a given vector $u \in \mathcal{X}$, we can define a diagonal operator $\text{Diag}(u) \in \mathcal{L}(\mathcal{X})$ to denote the diagonal operator whose diagonal entries are given by the entries of $u$.
That is, we have
\begin{equation}
    \text{Diag}(u)_{a, b} = \begin{cases}
        u_a & \text{if } a = b \\
        0 & \text{if } a \neq b
    \end{cases}
\end{equation}
\end{definition}

\begin{remark}{projection operators}\label{remark:projection_operators}
    Let $\Pi$ denote the set of projection operators MARK
\end{remark}

\begin{remark}{isometries}\label{remark:isometries}
    An operator $\mathcal{A} \in \mathcal{L}(\mathcal{X}, \mathcal{Y})$ is an isometry if it satisfies the condition
    \begin{equation}
        \mathcal{A}^{*} \mathcal{A} = I_{\mathcal{X}},
    \end{equation}
    which is equivalent to preserving the Euclidean norm of vectors in $\mathcal{X}$, i.e., $\|\mathcal{A}|\psi\rangle\| = \||\psi\rangle\|$ for all $|\psi\rangle \in \mathcal{X}$.
    % MARK: I used bra-ket notation above, but am realizing I am inconsistently jumping between the two.
\end{remark}


\begin{remark}{unitary operators}\label{remark:unitary_operators}
    Denote the set of isometries that map $\mathcal{X}$ to itself as the set of unitary operators $U(\mathcal{X})$.
    Every unitary operator $U \in U(\mathcal{X})$ is necessarily invertible and satisfies the normal condition
    \begin{equation}
        UU^{*} = U^{*}U = I_{\mathcal{X}}.
    \end{equation}
\end{remark}

\section{Permutation Invariance and Properties of Unitarily Invariant Measures}




%% END GOAL
\begin{theorem}{Schur-Weyl duality}\label{thm:schur_weyl_duality}
    Let $\mathcal{X}$ be a complex Euclidean space, $k$ be a positive integer, and $X \in \mathcal{L}(\mathcal{X}^{\otimes k})$ be an operator.
    Then
    \begin{equation}
        \text{Comm}(\mathcal{U}(\mathcal{X}), k) = \text{Span}\{W_{\pi} : \pi \in S_k \}
    \end{equation}
\end{theorem}
